\documentclass[12pt]{article}

\usepackage{sbc-template}
\usepackage{graphicx,url}
\usepackage[utf8]{inputenc}
\usepackage[brazil]{babel}

\sloppy

\title{Exemplo de Artigo no Formato SBC com Figura e Citações}

\author{Marcelo Ferreira Leda Filho -- 21953111\inst{1}}

\address{Instituto de Computação -- Universidade Federal do Amazonas (UFAM)\\
  Manaus -- AM -- Brasil
  \email{marcelo.leda@icomp.ufam.edu.br}
}

\begin{document}

\maketitle

\begin{abstract}
  Este artigo apresenta um exemplo simples de uso do template da SBC no Overleaf, incluindo inserção de figura, uso de citações bibliográficas e arquivo \texttt{.bib}.
\end{abstract}

\begin{resumo}
  Este documento mostra como estruturar um texto acadêmico em \LaTeX{} no formato da SBC, inserindo imagem, referências bibliográficas e citações no corpo do texto.
\end{resumo}

\section{Introdução}

O \LaTeX{} é amplamente utilizado na escrita de textos científicos e acadêmicos, especialmente em áreas como Computação, Matemática e Engenharia. Segundo \cite{knuth:84}, ele oferece alta qualidade tipográfica e grande controle sobre a estrutura do documento.

Além disso, a escrita acadêmica em Computação frequentemente exige a inserção de figuras, tabelas, algoritmos e referências bibliográficas. O template da SBC já oferece uma base adequada para esse tipo de produção.

\section{Inclusão de Figura}

A Figura \ref{fig:cameraman} apresenta uma imagem de exemplo incluída no documento. No seu projeto Overleaf, a imagem deve estar salva com o nome \texttt{cameraman.png}.

\begin{figure}[ht]
  \centering
  \includegraphics[width=0.5\textwidth]{cameraman.png}
  \caption{Imagem de exemplo incluída no artigo.}
  \label{fig:cameraman}
\end{figure}

Como discutido por \cite{lamport:94}, o uso correto de elementos gráficos melhora a clareza e a apresentação do texto acadêmico.

\section{Referências e Citações}

As citações bibliográficas podem ser feitas com o comando \verb|\cite|. Por exemplo, podemos citar um livro clássico de \LaTeX{} \cite{lamport:94} e também a obra de Knuth \cite{knuth:84}.

No template da SBC, as referências são inseridas com:

\begin{verbatim}
\bibliographystyle{sbc}
\bibliography{referencias}
\end{verbatim}

\noindent (\texttt{referencias} é o nome do teu arquivo \texttt{.bib})

Esse formato aparece no modelo da Overleaf da SBC.

\section{Conclusão}

Este exemplo apresentou a estrutura básica de um artigo no formato SBC com inclusão de figura, uso de citações e arquivo bibliográfico externo. Esse tipo de organização é útil para relatórios, artigos e trabalhos acadêmicos em Computação.

\bibliographystyle{sbc}
\bibliography{referencias}

\end{document}
